\documentclass[12pt]{article}
\usepackage[utf8]{inputenc}
\usepackage[T1]{fontenc}
\usepackage{hyperref}
\title{Wamesa Language Resources for K-12 Curriculum Development}
\author{Compiled by Claude Sonnet 5 (claude-sonnet-5) for Andr\'as Kornai}
\date{}
\begin{document}
\maketitle

\section{Introduction}
Wamesa (ISO 639-3 wad; Glottocode wand1267), often called Wandamen in the literature, is an Austronesian language spoken on the neck of the Doberai Peninsula (Bird's Head) in Indonesian New Guinea. Estimates run to 5,000--8,000 speakers. Indonesian is widely used for education and trade.

\section{Primary Educational Resources}
No school materials were found.

\section{Linguistic and Reference Materials}
\begin{itemize}
\item Wamesa language (Wikipedia): \url{https://en.wikipedia.org/wiki/Wamesa_language}
\item NSF grant page for Emily Gasser's project, which planned a reference grammar, a transcribed text corpus and a small Wamesa--English--Indonesian dictionary. I could not confirm that a published grammar exists: \url{https://grantome.com/grant/NSF/BCS-1153795}
\item PARADISEC collection EAG1, which includes a dictionary deposit dated 2019 (text file): \url{https://catalog.paradisec.org.au/repository/EAG1} and \url{https://catalog.paradisec.org.au/collections/EAG1/items/Dict}
\item Omniglot pages on Wamesa and Wandamen: \url{https://www.omniglot.com/writing/wamesa.htm} and \url{https://www.omniglot.com/writing/wandamen.htm}
\end{itemize}

\section{Religious and Community Resources}
Joshua Project reports Bible portions (1911--1994) but no New Testament: \url{https://joshuaproject.net/index.php/people_groups/15853}

\section{Digital and Audio Resources}
Global Recordings Network has audio-visual Bible lessons and short audio messages: \url{https://www.globalrecordings.net/en/language/1273} and \url{https://globalrecordings.net/en/language/wad}

\section{Curriculum Development Context}
Wamesa is still spoken daily, so a mother-tongue curriculum is plausible. Documentation exists (dictionary draft, recordings, grammar project) but no classroom materials were found. Contact with the documenting researcher and local communities would be the starting point.
\end{document}